\chapter{Prueba finita de las regiones de
\texorpdfstring{\(\pi\), \(e\) y \(\varphi\)}{pi, e y phi}}
\label{chap:regiones}
\index{regiones de pi@regiones de \(\pi\)}

La generación precedente ha fijado qué hacen clausura, propagación
y autoescala. Este bloque contiene la prueba finita que sostiene esa
narración: catálogo de regiones, multiplicidad, orientación, transporte y
prolongación. La evaluación arquimediana se aplica al final sobre el
catálogo ya construido y no reemplaza su genealogía.

\section{Proyección ternaria del catálogo}

Sea
\[
 T_{\rm cat}:\Omega_{\rm TPK}\longrightarrow\{0,\ldots,999\}^6
\]
la aplicación de catálogo construida en el \cref{chap:tpk}, y sea
\[
 \mathcal U_6:=\operatorname{im}(T_{\rm cat}).
\]
El censo exhaustivo da \(|\mathcal U_6|=468\). Para
\(U_6=(u_1,\ldots,u_6)\in\mathcal U_6\), definimos
\[
\Psi:\mathcal U_6\longrightarrow\mathbb F_3^6,
\qquad
\Psi(U_6)=((-u_1)\bmod3,\ldots,(-u_6)\bmod3).
\]
La selección interna del estado enriquecido de TPK, demostrada en el capítulo
precedente, fija las tres etiquetas
\[
w_\pi=010211,
\qquad
w_e=201101,
\qquad
w_\varphi=121200.
\]
Esta palabra ternaria es el valor de una proyección definida sobre el
catálogo, no sobre todo el cubo entero \(\{0,\ldots,999\}^6\). Esta
precisión de dominio es necesaria: la preimagen de una palabra en el cubo
entero ambiente sería enorme,
mientras que las fibras empleadas a continuación son fibras relativas a
\(\mathcal U_6\). A cada elemento de la fibra el catálogo adjunta
multiplicidad, firma y estado de retorno.

\section{La fibra de \texorpdfstring{\(\pi\)}{pi}: cinco regiones distintas}

La microfibra relativa al catálogo
\[
\mathscr R_\pi^{\rm micro}
=\{U\in\mathcal U_6:\Psi(U)=010211\}
\]
contiene cinco regiones \(U_6\) distintas:
\[
\begin{aligned}
U_\pi^\perp&=(501,614,498,169,272,272),\\
U_{\pi,1}^{++}&=(810,923,870,169,272,272),\\
U_{\pi,2}^{++}&=(810,983,810,169,272,272),\\
U_{\pi,3}^{++}&=(870,923,810,169,272,272),\\
U_\pi^{--}&=(870,923,810,418,674,521).
\end{aligned}
\]
Sus firmas multiplicativas son, respectivamente,
\[
555555,\quad339555,\quad393555,\quad933555,\quad933717,
\]
y sus multiplicidades
\[
432,\quad144,\quad144,\quad144,\quad144.
\]

\begin{figure}[p]
\centering
\resizebox{.96\textwidth}{!}{\input{figures/fig_cinco_regiones_pi}}
\caption{Microfibra de \(\pi\). La proyección a \(w_\pi\) conserva la
palabra y elimina la diferencia regional.}
\label{fig:cinco-pi}
\end{figure}

Bajo el alineamiento regional declarado con el sistema de incidencia, las
etiquetas de rol se distribuyen como
\[
5=1_\perp+3_{++}+1_{--};
\]
por multiplicidad,
\[
1008=432+4\cdot144.
\]
Las cuatro primeras regiones conservan la cola \(169|272|272\); la última es
la rama de retorno mutado. Los identificadores de catálogo no son canónicos,
porque cambian al reordenar las filas. Los vectores \(U_6\) sí lo son dentro
del catálogo fijado.

\begin{remark}[Fibra regional y cociente ternario]
El cociente
\(\mathscr R_\pi^{w_6}=\{010211\}\) tiene un elemento porque elimina los datos
regionales. La microfibra tiene cinco porque los conserva. Escribir «la región
de \(\pi\)» sólo es admisible si se especifica cuál de los dos tipos se está
usando.
\end{remark}

\section{Las fibras simples de \texorpdfstring{\(e\)}{e} y
\texorpdfstring{\(\varphi\)}{phi}}

La misma proyección que produce la microfibra quíntuple de \(\pi\)
determina dos fibras simples:
\[
 \Psi^{-1}(201101)=\{U_e\},
 \qquad
 \Psi^{-1}(121200)=\{U_\varphi\},
\]
donde
\[
 U_e=(850,963,890,149,252,212),
 \qquad
 U_\varphi=(830,943,830,109,252,252).
\]
Ambas regiones tienen multiplicidad \(144\), clase de orientación aditiva
\(\mathrm{ES}\), soporte multiplicativo de cardinal seis y tipo residual
\(A\). Sus firmas multiplicativas
son, respectivamente,
\[
 e_\times(U_e)=771339,
 \qquad
 e_\times(U_\varphi)=555933.
\]
La reducción directa de sus coordenadas es
\[
 U_e\bmod3=102202,
 \qquad
 U_\varphi\bmod3=212100;
\]
por la convención \(\Psi(U)=(-U)\bmod3\), estas dos palabras se convierten
en \(201101\) y \(121200\). Así, la diferencia entre la fibra de \(\pi\) y
las de \(e,\varphi\) ya está presente antes de toda evaluación decimal:
\(\pi\) posee cinco realizaciones regionales con multiplicidad total
\(1008\), mientras que \(e\) y \(\varphi\) poseen una realización cada una.

\section{Caracteres estructurales de las tres fibras}

Los datos anteriores distinguen tres comportamientos combinatorios. La fibra
de \(\pi\) registra clausura y orientación porque contiene ramas de retorno
diferentes; la de \(e\) queda caracterizada, después de aplicar la carta
arquimediana, por una complementariedad nonádica; y la de \(\varphi\) se
compara con el operador clásico de autoescala. Esta lectura se resume como
\[
\begin{array}{ccl}
\pi&:&\text{clausura, borde y orientación},\\
e&:&\text{propagación y crecimiento},\\
\varphi&:&\text{autoescala y proporción},\\
\alpha&:&\text{acoplamiento de las tres evaluaciones mediante el vector }K.
\end{array}
\]
La primera asignación es un invariante del catálogo regional. Las otras dos
utilizan los operadores arquimedianos que se introducen en las secciones
siguientes: la complementariedad de \(e\) y el radio espectral de
\(\varphi\). El cierre dodecafásico construye después el vector \(K\), la
frontera y la identidad conjunta que genera \(\alpha\); esta prueba finita no
introduce una calibración paralela de esos objetos.

Las doce tríadas arquimedianas que alimentan el cierre dodecafásico son
\[
\begin{aligned}
P={}&(141,592,653,589,793,238,462,643,383,279,502,884),\\
E={}&(718,281,828,459,045,235,360,287,471,352,662,497),\\
\Phi={}&(618,033,988,749,894,848,204,586,834,365,638,117).
\end{aligned}
\]
Estos tres vectores no son salidas de \(\Psi\): pertenecen al segundo nivel
de la construcción. El control finito que se especifica más abajo reproduce
sus cuatro primeras tríadas. Las doce tríadas completas son la salida de la
generación monodrómica del capítulo precedente y no se introducen como datos
de este control de cuatro pasos. En particular, \(P\) es la coordenada
decimal común de las cinco regiones de \(\pi\), pero no las identifica en el
espacio completo de estados.

\section{La región de \texorpdfstring{\(e\)}{e} y la complementariedad \texorpdfstring{\(1000-1\)}{1000-1}}

La primera ventana arquimediana asignada a \(w_e\) es
\[
718|281|828|459.
\]
Entre las \(243\) palabras de
\(\mathcal W_6:=\Psi(\mathcal U_6)\), y después de aplicar el control finito
de la carta arquimediana definido en la sección siguiente, \(w_e\) es la única
cuya primera pareja de tríadas suma \(999\). Además,
\[
718=729-11,
\qquad
281=271+10,
\qquad
718+281=999=1000-1,
\]
\[
828=3\cdot271+15,
\qquad
459=\sum R^\times.
\]
La caracterización finita es exacta una vez fijados el catálogo, los
operadores del TPK y la condición de complementariedad. Al aplicar el
predicado no se introduce otro valor decimal; esta unicidad es un control
interno de la generación ya construida y no una segunda derivación de \(e\).

El objeto examinado en esta sección es la órbita finita de cuatro
transiciones. La extensión a profundidad arbitraria y su compatibilidad con
la supervivencia no se infieren de esta ventana: pertenecen a la relación
EF--\(G_9\) y al teorema monodrómico del capítulo precedente.

\section{\texorpdfstring{\(\varphi\)}{phi}, Fibonacci y nivel cinco}

La proporción áurea es el radio espectral de
\[
N_\tau=\begin{pmatrix}0&1\\1&1\end{pmatrix},
\]
y el cociente continuo asociado satisface el punto fijo
\[
\varphi=1+\frac1\varphi.
\]
Estas son identidades clásicas. Una extensión posible considera el nivel
cinco de la correspondencia pentada--hexada, las clases
de Rogers--Ramanujan y las reglas de Fibonacci. Para convertir la convergencia
\(R(q)\to\varphi^{-1}\) en información HMT no basta evaluar cerca de
\(q=1\), porque el límite es universal. El observable debe ser diferencial o
espectral después de sustraer la asintótica dominante. La construcción de tal
observable y de su entrelazador con los productos modulares constituye una
hipótesis adicional a los resultados de esta sección.

\begin{remark}[Correspondencia pentada--hexada]
Una ruta concreta consiste en seleccionar mediante HMT un elemento de orden
cinco en la acción Witt, descomponer sus modos en las clases
\(\{\pm1\}\) y \(\{\pm2\}\pmod5\), y demostrar que las funciones
generatrices resultantes son los dos productos de Rogers--Ramanujan. Si el
operador de transferencia es \(N_\tau\), \(\varphi\) aparecerá como radio
espectral y la relación con anyones de Fibonacci será una representación, no
una analogía verbal.
\end{remark}

\section{Verificación finita de la contracción arquimediana sobre cilindros regionales}

El orden de tipos es
\[
\begin{aligned}
\text{admisibilidad}
&\to\text{semillas}
\to\text{palabra corta}
\to\text{calendario de continuación}\\
&\to\text{cilindro}
\to\text{supervivencia}
\to\text{evaluación real}.
\end{aligned}
\]
La cadena decimal aparece al final. Para hacer explícita la evaluación
finita, fijamos sobre \(\mathbb F_3^6\), con vectores fila y acción por la
derecha, las matrices
\[
L_0=
\begin{pmatrix}
2&2&2&1&2&1\\
2&1&2&2&1&1\\
1&0&0&1&0&2\\
0&0&1&1&1&0\\
2&2&2&0&2&1\\
2&1&2&1&0&0
\end{pmatrix},
\qquad
L_1=
\begin{pmatrix}
2&1&2&0&2&2\\
1&2&1&1&2&0\\
1&2&1&1&1&2\\
0&1&0&0&2&1\\
2&0&2&1&0&1\\
0&2&2&2&1&1
\end{pmatrix}.
\]
Estas matrices son operadores del estado enriquecido de TPK. No son parámetros que
se ajusten durante la evaluación: fijados el calendario y el estado inicial,
la lectura es determinista. Además, las seis transiciones publicadas de rango
completo caracterizan cada operador de manera única,
\[
L=X^{-1}Y\qquad\text{en }\mathbb F_3.
\]
Esta caracterización dual pertenece al estado enriquecido de TPK. Obtener de nuevo todos
sus campos desde la proyección reducida APP mínima sería un teorema de
realización distinto y no convierte \(L_0,L_1\) en datos introducidos a mano.

Para \(w=(d_0,\ldots,d_5)\in\mathbb F_3^6\) y un calendario
\(c=(c_1,c_2,c_3,c_4)\in\{0,1\}^4\), sean
\[
 u_0=w,
 \qquad
 u_j=u_{j-1}L_{c_j}\quad(1\leq j\leq4),
 \qquad
 W_c(w)=u_0\Vert u_1\Vert u_2\Vert u_3\Vert u_4.
\]
Si \(W_c(w)=(a_0,\ldots,a_{29})\), definimos mediante aritmética entera
\[
 N_c(w)=\sum_{k=0}^{29}a_k3^{29-k},
 \qquad
 \mathcal A_c(w)_j=
 \left\lfloor\frac{1000^jN_c(w)}{3^{30}}\right\rfloor\bmod1000,
 \quad1\leq j\leq4.
\]
La aplicación
\(\mathcal A_{0011}:\mathbb F_3^6\to\{0,\ldots,999\}^4\) es, por tanto,
la ventana finita de la carta arquimediana. Su dominio y codominio difieren de
los de \(\Psi\), y su definición abarca cuatro iteraciones. La prolongación a
profundidad arbitraria no se postula a partir de esta ventana: la proporciona
la conexión EF--\(G_9\) ya construida.

El censo exacto de las \(729\) palabras verifica que \(\mathcal A_{0011}\)
es inyectiva y que
\[
\begin{array}{c|c}
w&\mathcal A_{0011}(w)\\ \hline
010211&(141,592,653,589)\\
201101&(718,281,828,459)\\
121200&(618,033,988,749)
\end{array}
\]
son reconocimientos únicos. Entre los dieciséis calendarios binarios,
\(0011\) es el único que conserva simultáneamente estos tres
reconocimientos. Las \(144\) perturbaciones obtenidas al sustituir una entrada
de \(L_0\) o \(L_1\) por uno de los otros dos residuos destruyen el triple
reconocimiento. Los órdenes son
\[
\operatorname{ord}(L_0)=80,
\quad
\operatorname{ord}(L_1)=26,
\quad
\operatorname{ord}(L_0L_1)=242.
\]

La misma comprobación algebraica fija el estatuto lógico de este resultado.
Para \(L_0\), los seis vectores de entrada que proceden de las dos
primeras transiciones de \(w_\pi,w_e,w_\varphi\) tienen rango seis sobre
\(\mathbb F_3\); para \(L_1\) ocurre lo mismo con los seis vectores de las dos
transiciones siguientes. Por ello, las transiciones determinan cada matriz de
manera única. La unicidad, los órdenes y la rigidez frente a perturbaciones
son propiedades exactas del operador fijado. El cálculo directo distingue
esta caracterización dentro del estado completo de una eventual
reconstrucción adicional desde APP mínima.

La proyección observable de APP y el espacio completo de estados se distinguen
mediante la proyección de estado definida en el \cref{chap:tpk}; la carta
arquimediana es una aplicación posterior sobre el cociente ternario.

\begin{remark}[Catálogo y evaluación]
El resultado finito probado en esta sección tiene la forma
\[
\Omega_{\rm TPK}\xrightarrow{T_{\rm cat}}\mathcal U_6
\xrightarrow{\Psi}\mathcal W_6
\xrightarrow{\mathcal A_{0011}}\{0,\ldots,999\}^4.
\]
Las dos primeras flechas constituyen el catálogo; la tercera es el control
finito de la carta. Conservar esta separación impide atribuir al censo regional una
afirmación que sólo se ha verificado después de fijar \(L_0,L_1\).
\end{remark}
